% ============================
% --- BLOCK 1 (kept exactly) --
% ============================

\usepackage{amsmath,amsfonts,bm}

% Mark sections of captions for referring to divisions of figures
\newcommand{\figleft}{{\em (Left)}}
\newcommand{\figcenter}{{\em (Center)}}
\newcommand{\figright}{{\em (Right)}}
\newcommand{\figtop}{{\em (Top)}}
\newcommand{\figbottom}{{\em (Bottom)}}
\newcommand{\captiona}{{\em (a)}}
\newcommand{\captionb}{{\em (b)}}
\newcommand{\captionc}{{\em (c)}}
\newcommand{\captiond}{{\em (d)}}

% Highlight a newly defined term
\newcommand{\newterm}[1]{{\bf #1}}

% Figure reference, lower-case.
\def\figref#1{figure~\ref{#1}}
% Figure reference, capital.
\def\Figref#1{Figure~\ref{#1}}
\def\twofigref#1#2{figures \ref{#1} and \ref{#2}}
\def\quadfigref#1#2#3#4{figures \ref{#1}, \ref{#2}, \ref{#3} and \ref{#4}}
% Section reference
\def\secref#1{section~\ref{#1}}
\def\Secref#1{Section~\ref{#1}}
\def\twosecrefs#1#2{sections \ref{#1} and \ref{#2}}
\def\secrefs#1#2#3{sections \ref{#1}, \ref{#2} and \ref{#3}}
% Equation reference
\def\eqref#1{equation~\ref{#1}}
\def\Eqref#1{Equation~\ref{#1}}
\def\plaineqref#1{\ref{#1}}
% Chapter reference
\def\chapref#1{chapter~\ref{#1}}
\def\Chapref#1{Chapter~\ref{#1}}
\def\rangechapref#1#2{chapters\ref{#1}--\ref{#2}}
% Algorithm reference
\def\algref#1{algorithm~\ref{#1}}
\def\Algref#1{Algorithm~\ref{#1}}
\def\twoalgref#1#2{algorithms \ref{#1} and \ref{#2}}
\def\Twoalgref#1#2{Algorithms \ref{#1} and \ref{#2}}
% Part reference
\def\partref#1{part~\ref{#1}}
\def\Partref#1{Part~\ref{#1}}
\def\twopartref#1#2{parts \ref{#1} and \ref{#2}}

\def\eps{{\epsilon}}
\let\ab\allowbreak

% ============================
% --- BLOCK 2 MERGED BELOW ---
% ============================

\usepackage[a4paper,margin=0.8in]{geometry}

% Math & Symbols (remove duplicates: amsmath, amsfonts, bm)
\usepackage{amssymb, amsthm, mathtools}
\usepackage[most]{tcolorbox}
\usepackage{mathrsfs}
% needs texlive-science: \usepackage{algorithm}
% needs texlive-science: \usepackage{algpseudocode}
\usepackage{float}

% Colors & Links
\usepackage[dvipsnames]{xcolor}
\usepackage[colorlinks=true, linkcolor=blue, urlcolor=blue]{hyperref}

% Lists
\usepackage{enumitem}
\setlist{
    itemsep=0.6ex,
    topsep=1.6ex,
    partopsep=0.2ex,
    parsep=0.2ex
}

% Graphics
\usepackage{graphicx}
\usepackage{subcaption}

% Code Listings
\usepackage{listings}
\definecolor{codegray}{rgb}{0.95,0.95,0.95}
\definecolor{codepurple}{rgb}{0.58,0,0.82}
\definecolor{codegreen}{rgb}{0,0.6,0}

\lstdefinestyle{cppcode}{
    backgroundcolor=\color{codegray},
    commentstyle=\color{codegreen},
    keywordstyle=\color{magenta},
    numberstyle=\tiny\color{gray},
    stringstyle=\color{codepurple},
    basicstyle=\ttfamily\footnotesize,
    breaklines=true,
    numbers=left,
    numbersep=5pt,
    showspaces=false,
    showstringspaces=false,
    showtabs=false,
    tabsize=2
}
\lstset{style=cppcode}

\lstdefinestyle{pycode}{
    language=Python,
    basicstyle=\ttfamily\small,
    backgroundcolor=\color{gray!10},
    keywordstyle=\color{blue}\bfseries,
    commentstyle=\color{OliveGreen}\itshape,
    stringstyle=\color{red},
    numbers=left,
    numberstyle=\tiny\color{gray},
    frame=single,
    breaklines=true,
    tabsize=4,
    captionpos=b,
}

\lstnewenvironment{pycode}[1][]
{\lstset{style=pycode,#1}}
{}

\definecolor{darkpink}{RGB}{200,0,120}

\usepackage{tikz}
\usetikzlibrary{arrows.meta, positioning, shapes.geometric, calc, fit}
\usetikzlibrary{shapes.geometric}
\usetikzlibrary{arrows.meta}
\usetikzlibrary{positioning}
\usetikzlibrary{fit}
\usetikzlibrary{calc}
\usepackage{booktabs}
\usepackage{multirow}
\usepackage{array}

\usepackage{xparse}
\newcommand{\qtitle}[1]{\refstepcounter{section}\section*{\S\thesection: #1}
\addcontentsline{toc}{section}{\protect\numberline{\S\thesection: }#1}}

\newcounter{problem}
\NewDocumentEnvironment{problem}{o}
{\par\medskip\noindent\refstepcounter{problem}%
 \IfNoValueTF{#1}
   {\textbf{\textcolor{darkpink}{Problem \theproblem.}}~}
   {\textbf{\textcolor{darkpink}{Problem \theproblem~(#1).}}~}%
 \addcontentsline{toc}{section}{Problem \theproblem}}
{\par\medskip}

\NewDocumentEnvironment{solution}{}%
{\par\medskip\noindent\textbf{Solution.}\\}
{\par\medskip}

% Note environment
\newtcbtheorem[number within=section]{note}{Note}{
    enhanced,
    colback = black!10,
    colbacktitle = black!80,
    coltitle = black!5,
    boxrule = 0pt,
    frame hidden,
    fonttitle = \bfseries\sffamily,
    breakable,
    before skip = 2ex,
    after skip = 2ex
}{note}

\usepackage{verbatim}

\newcommand{\refnote}[1]{see \ref{#1}}

\newtheorem{theorem}{Theorem}
\newtheorem{definition}{Definition}

\usepackage{tocloft}

% Reduce vertical space between TOC lines
\setlength{\cftbeforesecskip}{8pt}   % default ~12pt
\setlength{\cftbeforesubsecskip}{4pt}
\setlength{\cftbeforesubsubsecskip}{4pt}
